\textbf{Theoretical considerations:}

The incident particle flux on a wall (i.e.\ a certain direction on a surface)
is:
\[ \Phi = m_0 \cdot \Omega
   = \frac{1}{6}\cdot m_0 \cdot n \cdot S \cdot \bar{v} \tag*{(3)} \]
\[ m_0 \cdot n = m_0 \cdot \frac{N}{a^3} = \frac{m_0 N}{a^3}
   = \frac{m}{V} = \rho, \]
where: $m_0$ is the mass of one molecule; $m$ -- mass of the gas in the cube;
$V$ -- volume of the cube; $\rho$ -- the density of the gas;
\[ \rho = \frac{\mu p}{RT}, \]
where $p$ -- the pressure of the gas in the cube. The relation (3) -- the mass
flux relation becomes:
\[ \Phi = \frac{1}{6}\cdot\rho\cdot S\cdot\bar{v}
   = \frac{1}{6}\cdot\frac{\mu p}{RT}\cdot S\cdot\sqrt{\frac{3RT}{\mu}}
   = \frac{1}{6}\cdot p\cdot S\cdot\sqrt{\frac{3\mu}{RT}}. \tag*{(4)} \]
\[ \Phi = \frac{1}{6}\cdot\rho\cdot S\cdot\bar{v}
   = \frac{1}{6}\cdot\frac{\mu p}{RT}\cdot S\cdot\sqrt{\frac{3RT}{\mu}}
   = \frac{1}{6}\cdot p\cdot S\cdot\sqrt{\frac{3\mu}{RT}}. \tag*{(5)} \]
% sic: relation (5) is a verbatim repeat of relation (4) in the source
\hfill(3 points)

\textbf{B. Reasoning}

Because the cylindrical can is not full of ice, in the empty part of it there
are saturated vapors, i.e.\ the mass flux of the molecule which sublimate is
equal with mass flux of gass % sic: "gass" in source
which transform into ice. Thus the pressure in the can is the saturated vapor
pressure $p_{\mathrm{s}}$ and it has the corresponding maximum density
$\rho_{\mathrm{s}}$. See figure 6.2.
% FIGURE NEEDED: Fig. 6.2 -- two cylindrical cans: (a) closed can with ice ("gheata") at the bottom and saturated vapors ("vapori saturati") at pressure p_s above, opposite fluxes Phi_1 (up) and Phi_2 (down); (b) opened can with pressure p_s/2 and only the flux Phi_1 (2014 TH solutions, p. 10)

\begin{align*}
  \Phi_1 &= \Phi_{\mathrm{sublimation}}
    = \frac{1}{6}\cdot\rho_{\mathrm{s}}\cdot S\cdot\bar{v}
    = \frac{1}{6}\cdot p_{\mathrm{s}}\cdot S\cdot\sqrt{\frac{3\mu}{RT}}; \\
  \Phi_2 &= \Phi_{\mathrm{solidification}}
    = \frac{1}{6}\cdot\rho_{\mathrm{s}}\cdot S\cdot\bar{v}
    = \frac{1}{6}\cdot p_{\mathrm{s}}\cdot S\cdot\sqrt{\frac{3\mu}{RT}}; \\
  \Phi_1 &= \Phi_2.
\end{align*}

After the can was opened, there be no molecules which sublimate thus the mass
flux of the molecules which gather the ice become null. So the pressure
becomes $\left(p_{\mathrm{s}}/2\right)$.

Thus the force acting on the astronaut will be

\textbf{C. Calculations according to the reasoning}, and/or as support for
reasoning \hfill(2 points)
\[ F = \frac{p_{\mathrm{s}}}{2}\cdot S, \]
Opening the can the astronaut will be accelearated % sic: "accelearated" in source
with:
\[ a = \frac{F}{M} = \frac{p_{\mathrm{s}}\cdot S}{2M}
   = \frac{550\,\mathrm{N\,m^{-2}} \cdot 30\cdot10^{-4}\,\mathrm{m^2}}
          {2\cdot10^{2}\,\mathrm{kg}}
   = 0{,}00825\,\mathrm{m\,s^{-2}}. \]

The total time of the acceleration movement will be the total time of ice
sublimation:
\[ \tau = \frac{m}{\Phi_1}
   = \frac{m}{\dfrac{1}{6}\cdot p_{\mathrm{s}}\cdot S
     \cdot\sqrt{\dfrac{3\mu}{RT}}}
   = \frac{6m}{p_{\mathrm{s}}\cdot S}\cdot\sqrt{\frac{RT}{3\mu}}
   \approx 150\,\mathrm{s}. \]

\textbf{D. Correct result}

The travel distance in this time will be:
\[ L = \frac{a\tau^2}{2}
   = \frac{0{,}00825\,\mathrm{m\,s^{-2}} \cdot 225\cdot10^{2}\,\mathrm{s^2}}{2}
   \approx 93\,\mathrm{m}, \]
The astronaut could arrive safely in at the shuttle if he didn't lose to much
time by solving the problem. % sic: "in at the shuttle", "to much" in source
